\clearpage\section{275+ Follow-On Research Questions}
The 60 papers generated 275+ follow-on research questions (five per paper for the original 45, per the rubric; Wave~3 biology and Wave~4 PDE questions are included in-line with their respective paper subsections in \S\S46--60). We group the original 225 questions by domain bucket below. Each block is headed by the paper; questions are numbered within paper.
\subsection*{Astrophysics/Cosmology}
\addcontentsline{toc}{subsection}{Astrophysics/Cosmology}
\paragraph{Approximating Photo-z PDFs for Large Surveys \textnormal{\small(OSTI \texttt{1461824})}}
\begin{enumerate}[leftmargin=1.6em,itemsep=0.2em,topsep=0.2em,label=Q\arabic*.]
\item Generate BPZ-like mocks from a real photometric catalog (e.g., COSMOS2020 or HSC-SSP) and re-run the KLD benchmark --- does the quantile dominance hold when the PDF shape distribution comes from real data
\item At fixed byte budget (not N\_f), which format wins Compute exact byte cost including metaparameters and float precision (float16 vs float32) --- does quantiles' advantage shrink when samples use int16 indices
\item Train a neural PDF compressor (VAE or normalizing flow) on the same mock catalogs --- does a learned 10-dim latent beat the best classical format at matched byte budget
\item How does quantile reconstruction degrade under realistic tail truncation (z quantile levels clipped to [0.01, 0.99]) Sweep clipping thresholds and identify the regime where samples format surpasses quantiles.
\item Propagate each compressed format through a 2-point cosmic-shear cosmological inference --- does the quantile format's lower KLD translate to tighter $\Omega$\_m, $\sigma$ contours, and by how many percent
\end{enumerate}
\paragraph{Latent Stochastic Differential Equations for Modeling Quasar Variability and Inferring Black Ho \textnormal{\small(OSTI \texttt{2396968})}}
\begin{enumerate}[leftmargin=1.6em,itemsep=0.2em,topsep=0.2em,label=Q\arabic*.]
\item How does the latent SDE's parameter recovery degrade as photometric noise scales to 10x LSST requirement
\item Can replacing the GRU encoder with a Transformer (time-series foundation model like TimesFM) improve parameter recovery for the weak-signal parameters (a, lambda\_Edd)
\item What is the minimum survey duration (vs LSST 10yr) needed to recover log(M\_BH) with 0.1 dex precision
\item Does adding broad-line region emission (reverberation mapping light curves) as an auxiliary input tighten black hole mass constraints
\item How robust is the latent SDE posterior calibration under model mismatch (e.g., the accretion disk is not a simple thin-disk lamppost)
\end{enumerate}
\paragraph{Cosmic Reionization On Computers: Properties of the Post-Reionization IGM \textnormal{\small(OSTI \texttt{1275503})}}
\begin{enumerate}[leftmargin=1.6em,itemsep=0.2em,topsep=0.2em,label=Q\arabic*.]
\item Can a learned correction to the FGPA lognormal tail (e.g. a small neural network trained to map 1D lognormal -\textgreater{} CROC flux PDF) close the z\textgreater{}5.5 gap while retaining the 10\textasciicircum{}4-fold speedup
\item Inject a toy proximity-zone population (Poisson-distributed bright QSOs with PDF-convolved ionization bubbles) and test whether the missing wide-transmission-peak population is explained.
\item Re-calibrate the UV-background fluctuation field (sigma\_Gamma, ell\_mfp) by fitting simultaneously to flux PDF + dark-gap statistics at z\textless{}5.5 and check whether the fit is unique or degenerate.
\item Extend the semi-analytic model down to z=4 and up to z=7 and compare against the paper's full redshift range --- does the disagreement scale with the mean transmitted flux
\item Use the DC-mode test to derive an effective b(z,delta) bias of the Lyman-alpha forest in the reionization epoch and compare to theoretical predictions from the patchy-reionization literature.
\end{enumerate}
\subsection*{Bio/Bioinformatics}
\addcontentsline{toc}{subsection}{Bio/Bioinformatics}
\paragraph{Clustering huge protein sequence sets in linear time \textnormal{\small(OSTI \texttt{1624105})}}
\begin{enumerate}[leftmargin=1.6em,itemsep=0.2em,topsep=0.2em,label=Q\arabic*.]
\item How does Linclust's clustering quality (F1 vs curated reference families) degrade as intra-family identity is pushed from 50\% down to 30\% on a controlled synthetic benchmark
\item Can the bottom-m sketch be replaced with a learned MinHash via a small Transformer embedding, and does that improve remote-homolog recall at matched runtime
\item What is the GPU speedup achievable by porting the bucketed center-assignment stage to CUDA (hash-table on device memory), and does the memory bandwidth bound dominate
\item How sensitive is the linear-time scaling exponent to the k-mer length k and the sketch width m, and is there a regime where the method silently becomes O(N log N)
\item When inputs are metagenome-scale (\textgreater{}5e9 sequences with heavy sequence redundancy), does Linclust's single-pass center selection still produce stable cluster boundaries under random input shuffling
\end{enumerate}
\paragraph{Updated Virophage Taxonomy and Distinction from Polinton-like Viruses \textnormal{\small(OSTI \texttt{2475938})}}
\begin{enumerate}[leftmargin=1.6em,itemsep=0.2em,topsep=0.2em,label=Q\arabic*.]
\item Does adding the 32 newly published metagenome-assembled virophage genomes (2023-2024 IMG/VR release) shift the Mavirus-basal topology, or are the paper's conclusions robust
\item Can a protein language model (ESM-2 or ProtTrans) embedding of MCP and ATPase sequences distinguish virophages from PLVs/adintoviruses without HMM profiles, and at what sequence-identity floor does HMMER fail while LM embeddings still succeed
\item What is the minimum marker-gene subset (MCP alone vs 2/3/4 markers) needed to recover the paper's family-level topology at \textgreater{}= 95 UFBoot support
\item Under Bayesian concordance factor analysis, do the four marker trees support a single evolutionary history or reveal reticulation (horizontal transfer) at the Sputnik/Mavirus split
\item When reference-free clustering (vConTACT3 or mash-based) is applied to the paper's genome set, does it produce the same family boundaries, or does it over-/under-split relative to the HMM + phylogeny approach
\end{enumerate}
\paragraph{Supervised extraction of near-complete genomes from metagenomic samples: A new service in PATRI \textnormal{\small(OSTI \texttt{2469515})}}
\begin{enumerate}[leftmargin=1.6em,itemsep=0.2em,topsep=0.2em,label=Q\arabic*.]
\item Can a lightweight pheS-anchoring substitute be built from BV-BRC's public REST API (pull reference genomes on demand) to replicate the paper's supervised pipeline without requiring a local 193,980-genome install
\item On the same synthetic community, how does MetaBAT2 compare to modern contrastive-learning binners (SemiBin2, VAMB) in HQ-bin count and what does the paper's supervised pipeline need to beat them
\item At what point in the abundance-ratio skew does the differential-coverage advantage break down (e.g., flat-abundance community), and does the supervised pipeline's pheS anchor retain performance where coverage fails
\item Can a protein language model (ESM-2) replace the 1,300-role random-forest EvalCon consistency scorer at matched accuracy, and does it generalise to novel lineages without retraining
\item When the co-assembly is replaced with long-read (ONT or PacBio) single-sample assembly on the same synthetic community, does MetaBAT2 still produce 5/5 HQ bins or does the loss of differential-coverage information require the supervised pipeline to recover them
\end{enumerate}
\subsection*{CS/Graph Algorithms}
\addcontentsline{toc}{subsection}{CS/Graph Algorithms}
\paragraph{Mathematical Foundations of the GraphBLAS \textnormal{\small(OSTI \texttt{1379592})}}
\begin{enumerate}[leftmargin=1.6em,itemsep=0.2em,topsep=0.2em,label=Q\arabic*.]
\item Benchmark our Python GraphBLAS vs SuiteSparse:GraphBLAS on SNAP graphs (LiveJournal, Orkut, Friendster) for BFS, SSSP, triangle count --- what is the Python slowdown factor, and which primitives dominate
\item Express three recent graph algorithms (GNN message-passing, graph spectral clustering, personalized PageRank with teleport) as GraphBLAS primitives --- is the semiring abstraction sufficient, or what extensions are needed
\item Design a semiring for uncertainty propagation (value + variance pairs) and express variance-aware BFS and SSSP --- does the algebraic structure (associativity, distributivity) still hold, and does this yield provably-correct uncertainty quantification
\item Compare GraphBLAS-expressed triangle counting against vertex-centric (Pregel) and matrix-free (SuiteSparse LAGraph) implementations on a 10-edge graph --- at what scale does the matrix formulation beat vertex-centric, and where does each paradigm break
\item Formally verify (e.g., in Coq or Lean) that our 6 graph algorithms are correct for any abstract semiring satisfying the stated axioms --- this would provide the first machine-checked proof of GraphBLAS algorithm correctness.
\end{enumerate}
\paragraph{ScaWL: Scaling k-WL (Weisfeiler-Lehman) Algorithms in Distributed-Memory \textnormal{\small(OSTI \texttt{2587225})}}
\begin{enumerate}[leftmargin=1.6em,itemsep=0.2em,topsep=0.2em,label=Q\arabic*.]
\item Port the 2-WL implementation to C++/OpenMP on a 20-core node and reproduce the paper's 16x speedup at p=20 on the same synthetic graph --- quantify exact Python overhead.
\item Scale to 3-WL on small graphs (n=40-80) using the rank-reduction storage trick and compare empirical memory to the paper's O(k|V|\textasciicircum{}k) claim.
\item Benchmark 2-WL / 3-WL / 4-WL expressivity on the standard BREC graph-isomorphism benchmark --- how many more pairs does k=3 vs k=2 distinguish in practice
\item Replace deterministic canonicalization with approximate / sketched hashing (e.g. MinHash on neighborhood multisets) and measure speed vs correctness tradeoff on large sparse graphs.
\item Evaluate ScaWL as a subroutine inside a GNN (Provably Powerful Graph Networks) --- does distributed k-WL enable larger-graph training and does this improve downstream accuracy on OGB benchmarks
\end{enumerate}
\subsection*{Computational Chemistry}
\addcontentsline{toc}{subsection}{Computational Chemistry}
\paragraph{Markov State Models from Short Non-Equilibrium Trajectories via Observable Operator Models \textnormal{\small(OSTI \texttt{1565592-MSM-Hempel})}}
\begin{enumerate}[leftmargin=1.6em,itemsep=0.2em,topsep=0.2em,label=Q\arabic*.]
\item Is the 44\% t gap (2020 vs 1400 ps) attributable to force field, water model, or integrator --- test by running identical short-trajectory protocol with CHARMM36m and GROMOS54a7 and comparing long-reference t.
\item How does OOM correction scale to $\geq$3 metastable states Apply to chignolin or villin headpiece where $\geq$3 slow modes exist and compare OOM vs direct MSM recovery of t, t, t.
\item Does the 'both-basin-coverage' clustering requirement we discovered translate into a principled seeding algorithm (e.g., farthest-point after long-ref preseeding) that recovers OOM accuracy without requiring a long reference
\item Benchmark OOM vs Koopman-based VAMPnets on identical short-trajectory ADP data --- which method recovers t with fewer trajectories, and how does the crossover depend on trajectory length
\item Quantify the breaking point: at what short-trajectory length T\_min does OOM no longer correct the MSM bias (error \textgreater{} 20\% of long-ref t) Sweep T $\in$ \{1,5,10,20,50\} ps at fixed total data.
\end{enumerate}
\subsection*{Control Theory}
\addcontentsline{toc}{subsection}{Control Theory}
\paragraph{Motion Tomography via Occupation Kernels \textnormal{\small(OSTI \texttt{1842593})}}
\begin{enumerate}[leftmargin=1.6em,itemsep=0.2em,topsep=0.2em,label=Q\arabic*.]
\item Download a public oceanographic glider dataset (e.g., IOOS Glider DAC) and apply Algorithm 1 to real endpoint displacements --- does kernel width $\mu$=0.5 remain optimal on physical-scale flows, and what is the reconstruction error vs HYCOM model ground truth
\item Add Gaussian observation noise $\sigma$ to endpoint displacements and characterize algorithm breakdown --- at what $\sigma$/displacement ratio does the reconstruction error exceed 20\%, and does Tikhonov regularization extend the tolerable noise range
\item Replace the occupation-kernel RKHS with a divergence-free neural-network parameterization --- does the learned approach beat the kernel method on the Gaussian-mixture flow at N=25, and how does it scale to 4D (spatial+time) flows
\item Extend Algorithm 1 to time-varying flows F(x,t) via space-time kernels K((x,t),(y,s)) --- does it recover a known oscillating-vortex test case, and what is the effective temporal resolution as a function of trajectory density
\item For the Gaussian-mixture flow, characterize identifiability: with fixed trajectory endpoints, what subspace of F is unrecoverable, and does this null space correspond to a closed-form RKHS orthogonal complement
\end{enumerate}
\subsection*{Engineering/Physics}
\addcontentsline{toc}{subsection}{Engineering/Physics}
\paragraph{Numerical study of the ignition behavior of a post-discharge kernel in a turbulent stratified c \textnormal{\small(OSTI \texttt{1559043})}}
\begin{enumerate}[leftmargin=1.6em,itemsep=0.2em,topsep=0.2em,label=Q\arabic*.]
\item When the simulated window is extended to t = 500 us with AMR active (refinement around OH \textgreater{} 1e-3 isocontour), does the single-realisation case at phi = 1.0 transition to self-sustained propagation, and what is the measured ignition delay
\item How does the monotonic OH-vs-phi ordering change if the kernel radical composition (Y\_OH = 0.01, Y\_O = 0.01, Y\_H = 0.001) is replaced with an equilibrium-plasma mixture at 3300 K, and does ignition probability become phi-insensitive
\item Can a low-cost analytic surrogate (kernel heat release + local Da number) predict the paper's ignition-probability curve from 2-3 PeleC realisations plus a calibrated stochastic model, reducing the required ensemble size by 5x
\item What is the sensitivity of the phi-ordering to the DRM19 -\textgreater{} GRI-3.0 chemistry upgrade, and does the additional NOx chemistry change the OH pool saturation plateau
\item Under cross-flow velocity scaling (10-40 m/s), does the critical phi for ignition shift monotonically with the local Da number, and can the paper's single-velocity result be collapsed onto a Da-based master curve
\end{enumerate}
\subsection*{Machine Learning}
\addcontentsline{toc}{subsection}{Machine Learning}
\paragraph{Physics-based hybrid machine learning for critical heat flux prediction with uncertainty quanti \textnormal{\small(OSTI \texttt{2571909})}}
\begin{enumerate}[leftmargin=1.6em,itemsep=0.2em,topsep=0.2em,label=Q\arabic*.]
\item Does the hybrid advantage under data scarcity persist if the base correlation is deliberately mis-specified (e.g., Biasi used outside its valid L/D range) and how quickly does performance degrade
\item For the 9-point regime, can transfer learning from a model pretrained on the synthetic surrogate produce calibrated uncertainty on the real NRC data with \textless{}50 fine-tuning points
\item How does epistemic-uncertainty miscalibration area scale with training set size for the three UQ backbones (DNN ensemble vs BNN vs DGP) at matched parameter count
\item Can symbolic regression over the learned residual recover an interpretable correction term to Biasi, and does that term resemble known Groeneveld corrections in the rod-bundle regime
\item When the input-feature set is reduced to operating conditions only (drop geometry), how much of the hybrid advantage disappears, quantifying the contribution of each feature family
\end{enumerate}
\paragraph{A Portfolio Approach to Massively Parallel Bayesian Optimization \textnormal{\small(OSTI \texttt{2571540})}}
\begin{enumerate}[leftmargin=1.6em,itemsep=0.2em,topsep=0.2em,label=Q\arabic*.]
\item Reproduce the noisy Branin/Hartmann6 benchmarks in Python with heteroscedastic GP --- does qHSRI's automatic replication allocation match the paper's claimed advantage over qAEI
\item Why does BoTorch's qEI underperform so dramatically (worse than RS) Instrument the joint-acquisition optimization and diagnose: is it L-BFGS local minima, the raw-samples heuristic, or kernel conditioning
\item Test qHSRI at q=100 on a 50D Ackley benchmark to verify the paper's O(1) batch-cost claim --- does wall-clock per iteration stay flat vs q, and does regret scale as $\sqrt{}$(log q)
\item Benchmark qHSRI against modern batch methods released after 2021 (TurBO, LaMBO, BORE) on the same 4 noiseless benchmarks with identical n, q, seeds --- where does the portfolio approach sit in the 2026 landscape
\item Combine Thompson sampling candidates (our surprisingly strong qTS) with Sharpe-ratio selection in a hybrid qTS-HSRI --- does the combined method beat both on Hartmann6-12D
\end{enumerate}
\paragraph{Divide and Conquer: Learning Chaotic Dynamical Systems with Multi-Step Penalty Neural ODEs (MP- \textnormal{\small(OSTI \texttt{3003857})}}
\begin{enumerate}[leftmargin=1.6em,itemsep=0.2em,topsep=0.2em,label=Q\arabic*.]
\item Rerun the gradient explosion demo at T=20 and 2000 steps (matching paper) on a single A100 / Spark to confirm the 10\textasciicircum{}8x gap --- does MP's advantage scale with T as exp(lambda\_1 * T)
\item Implement Linot-2023's stabilized NODE architecture with dilated convolutions (rates 1,2,3,4,3,2,1) and retest the KS joint-PDF; does this rescue the KS training collapse seen in the current replication
\item Compare MP-NODE against LSS (least-squares shadowing) on Lorenz-63 for exact gradient accuracy at fixed compute --- is MP's O(m+n) cost vs LSS's O(m*n\textasciicircum{}3) the only justification or is MP also more accurate
\item Study the mu schedule automatically: can we derive an adaptive mu update rule (e.g. based on the penalty/GT loss ratio) that eliminates the 'problem-dependent empirical' schedule the paper describes
\item Apply MP-NODE to a shallow-water / Lorenz-96 / Kuramoto system and test whether a single mu schedule transfers across problems or truly requires per-system tuning.
\end{enumerate}
\subsection*{Materials/Condensed Matter}
\addcontentsline{toc}{subsection}{Materials/Condensed Matter}
\paragraph{Quantitative Relationship between Polarization Differences and the Zone-Averaged Shift Photocur \textnormal{\small(OSTI \texttt{1523841})}}
\begin{enumerate}[leftmargin=1.6em,itemsep=0.2em,topsep=0.2em,label=Q\arabic*.]
\item Extend the numerical verification to a 3-band or 4-band tight-binding model (e.g. SSH+onsite, or a minimal GeS-like model) and check whether the identity generalizes via the Wilson-loop formulation of Eq. 17.
\item Compute the shift-current conductivity sigma\textasciicircum{}shift(omega) in the Rice-Mele model and verify the frequency integral equals (pi*e\textasciicircum{}3/hbar\textasciicircum{}2)*(P\_c - P\_v) modulo the winding contribution.
\item Apply the identity as a DFT+Wannier post-processing tool for a material like GeS or monolayer WS2: compute P\_c-P\_v from Wannier centers and compare to the direct frequency-integrated sigma\textasciicircum{}shift.
\item Study how the winding number W\_cv changes under SHG-relevant perturbations (strain, electric field) and predict a shift-current discontinuity as a material tuning knob.
\item Generalize to interacting systems: does the identity survive a Hartree-Fock or GW correction when P\_c-P\_v is evaluated with quasiparticle wavefunctions Test on a Hubbard-Rice-Mele numerical example.
\end{enumerate}
\paragraph{Chiral Spin Order in Kondo-Heisenberg Systems \textnormal{\small(OSTI \texttt{1412756})}}
\begin{enumerate}[leftmargin=1.6em,itemsep=0.2em,topsep=0.2em,label=Q\arabic*.]
\item Run a classical Monte Carlo on the Ising-alpha effective model (on a 2D lattice with coarse-grained coherence volumes) to numerically extract beta and test the claimed 2D Ising universality.
\item Generalize the mean-field ansatz to triangular and kagome Kondo-Heisenberg lattices and compute how the CSL dome shifts --- does kagome frustration widen the CSL window between J\_c and J\_AFM
\item Incorporate itinerant-electron Berry curvature in the CSL phase and predict the anomalous Hall conductivity sigma\_xy(T) below T\_c; compare to measured sigma\_xy in candidate chiral-spin materials.
\item Relax the nested-FS assumption (tilde\_J\_H(Q)=0) and scan filling to map how the CSL dome deforms away from perfect nesting; identify the filling window where CSL survives.
\item Couple the mean-field electronic structure to phonons and ask whether the CSL onset couples to a lattice distortion (magnetoelastic signature) observable in Bragg-peak splitting.
\end{enumerate}
\paragraph{Physics and Chemistry from Parsimonious Representations: Image Analysis via Invariant Variation \textnormal{\small(OSTI \texttt{2439897})}}
\begin{enumerate}[leftmargin=1.6em,itemsep=0.2em,topsep=0.2em,label=Q\arabic*.]
\item On a synthetic dataset with two coupled physical factors (a AND orientation-angle-of-a-strain-tensor), can the rVAE cleanly disentangle both into separate content dimensions, or does it still recruit only one
\item Benchmark the rVAE against a modern equivariant neural-network baseline (e.g. escnn SE(2)-equivariant CNN autoencoder) on the same hexagonal-lattice task --- does explicit equivariance beat the implicit-via-pose-head approach
\item Study the posterior-collapse failure mode systematically: how does the content-to-physics |r| depend on beta schedule shape (linear, exponential, KL-annealing frontend)
\item Apply the rVAE to a real open STEM dataset (e.g. the EMPIAD or NCEM public STEM datasets) and quantify whether the content latent recovers known crystallographic parameters.
\item Extend the pose head to non-trivial symmetry groups (C\_4, C\_6, or reflections) and measure the disentanglement gain on datasets with those specific symmetry groups vs generic SE(2).
\end{enumerate}
\paragraph{Deep Learning of Atomically Resolved Scanning Transmission Electron Microscopy Images: Chemical \textnormal{\small(OSTI \texttt{1427646})}}
\begin{enumerate}[leftmargin=1.6em,itemsep=0.2em,topsep=0.2em,label=Q\arabic*.]
\item When the synthetic training set is replaced with a small-angle-multislice abTEM simulation (3-4 slices, \textasciitilde{}10x cheaper than full Prismatic), does pixel accuracy hold and does generalisation to real STEM frames improve measurably
\item How does the U-Net's per-class F1 degrade on real experimental HAADF-STEM frames of SrTiO3/LSMO as dose is reduced from 1e5 to 1e3 e-/A\textasciicircum{}2, and at what dose does Sr vs La/Sr confusion become the dominant error
\item Can a self-supervised pretraining stage (masked autoencoder on unlabelled STEM frames) close the synthetic-to-real domain gap with \textless{}5\% real-frame labels
\item Does replacing the U-Net with a Vision Transformer + UperNet head preserve sub-pixel detection precision while providing calibrated classification uncertainty per atom column
\item When the model is applied across a temporal sequence of experimental STEM frames, does a simple Kalman smoother on per-column class probabilities recover transformation trajectories (antisite \textless{}-\textgreater{} vacancy) with known drift statistics
\end{enumerate}
\paragraph{Electronic and Optical Properties of Two-Dimensional GaN from First-Principles \textnormal{\small(OSTI \texttt{1484740})}}
\begin{enumerate}[leftmargin=1.6em,itemsep=0.2em,topsep=0.2em,label=Q\arabic*.]
\item Run GW@PBE on the existing Quantum ESPRESSO starting point using BerkeleyGW on a single Spark / small HPC slice --- how close is G0W0 monolayer gap to the paper's 6.32 eV
\item Redo the bilayer structural relaxation with tighter force tolerance (\textless{}1e-4 Ry/Bohr) and a larger vacuum + dipole correction; does the DFT gap recover toward 1.3 eV
\item Compute the biaxial strain dependence of the GW-corrected band gap (via scissor-proxy: assume GW correction is strain-independent) and compare to strain-dependent optical gap in GaN/AlN heterostructure experiments.
\item Replace H passivation with F or Cl and track how the band gap, polarization field, and exciton binding change --- identify a passivation that turns the bilayer direct-gap.
\item Use a machine-learned xc functional (e.g. SCAN, r2SCAN, or DM21) as a PBE alternative and quantify the DFT-functional uncertainty envelope on the 2D GaN gap before GW corrections.
\end{enumerate}
\paragraph{Generative AI-Driven Accelerated Discovery of Passivation Molecules for Perovskite Solar Cells \textnormal{\small(OSTI \texttt{PVMol-Gen-Fajar2026})}}
\begin{enumerate}[leftmargin=1.6em,itemsep=0.2em,topsep=0.2em,label=Q\arabic*.]
\item Does matching the paper's exact SMILES-X version + tokenizer close the F1 gap from 0.656 to 0.80, or is the gap due to training-data preprocessing (canonicalization, augmentation multiplier)
\item Run xTB on the 10 final candidates from both SMILES and SELFIES pipelines --- do the E\_gap/dipole distributions overlap with the paper's reported 10 molecules
\item Does training-set contamination from predicted class-1 relabeling (cycles 2-3) inflate apparent classifier accuracy Quantify by holding out a fixed external test set and tracking its F1 across cycles.
\item Substitute the GPT-2 generator with a chemistry-tuned foundation model (e.g., ChemGPT, MolGPT) --- does the class-1 rate at cycle 1 jump from 83\% to \textgreater{}95\%, and do top-10 clusters converge to similar motifs
\item Benchmark downstream perovskite figure-of-merit predictors (DFT binding energy on Cs/FA surfaces) on the 20 candidates (10 SMILES + 10 SELFIES) to test whether the representation choice biases toward different chemistries.
\end{enumerate}
\paragraph{Effect of Single Atom Platinum (Pt) Doping and Facet Dependent on the Electronic Structure and  \textnormal{\small(OSTI \texttt{1981773})}}
\begin{enumerate}[leftmargin=1.6em,itemsep=0.2em,topsep=0.2em,label=Q\arabic*.]
\item What is the facet-dependence of the Pt-induced band-gap narrowing when computed on (010) and (111) surfaces with matched slab thickness, and does (001) remain the most photoactive as the paper claims
\item Does a spin-polarised calculation of the Pt-doped (001) slab reveal a local magnetic moment on Pt, and does that shift the mid-gap states relative to the non-spin-polarised result
\item How does the gap narrowing evolve with Pt coverage (1/8, 1/4, 1/2, 1 ML) on the (001) surface, and at what coverage does the system become metallic
\item Using HSE06 or G0W0 on the relaxed Pt-doped slab, does the quantitative gap narrowing (absolute gap and CBM/VBM shifts) survive, or is it a PBE self-interaction artefact
\item Can a climbing-image NEB calculation of water-splitting intermediates (H* / OH*) on the Pt-doped (001) surface connect the electronic-structure change to a measurable overpotential difference vs the clean surface
\end{enumerate}
\subsection*{Mathematics}
\addcontentsline{toc}{subsection}{Mathematics}
\paragraph{Integer Sequences from Configurations in the Hausdorff Metric Geometry \textnormal{\small(OSTI \texttt{1997354})}}
\begin{enumerate}[leftmargin=1.6em,itemsep=0.2em,topsep=0.2em,label=Q\arabic*.]
\item Extend the exhaustive graph enumeration to K, and K, to resolve the 35 open integers in [1,100] --- which remain non-achievable
\item Is there a polynomial-time certificate for non-achievability, or does deciding achievability require exponential search Empirically test with SAT encoding on integers up to 200.
\item Compute the asymptotic density of achievable integers in [1,N] as N$\rightarrow$$\infty$ --- does it tend to 1, to a constant \textless{} 1, or to 0
\item Generalize the bijection to tripartite graph edge covers --- what new integer sequences arise, and do they correspond to existing OEIS entries
\item Which formula family (from Theorems 10--12) is tightest --- i.e., attains the minimum m+n realization for a given integer Produce an atlas mapping each achievable integer to its cheapest bipartite realization.
\end{enumerate}
\subsection*{Nuclear}
\addcontentsline{toc}{subsection}{Nuclear}
\paragraph{Molten Salt Reactor Neutronics and Fuel Cycle Modeling and Simulation with SCALE (MSRE Depletio \textnormal{\small(OSTI \texttt{2217719})}}
\begin{enumerate}[leftmargin=1.6em,itemsep=0.2em,topsep=0.2em,label=Q\arabic*.]
\item How sensitive is the 375-day reactivity swing to the online Xe-135 removal rate --- can we quantify dk/d(lambda\_Xe) and compare to MSRE operational data
\item Extend depletion to 900 days with periodic U-235 additions matching the MSRE operational log and compare drift in k\_eff to the paper's full-history curve.
\item Swap the fuel to a thorium-bearing MSR composition (LiF-BeF2-ThF4-UF4) and compute the breeding ratio with and without online protactinium separation.
\item Perform a Monte-Carlo-vs-deterministic cross-validation by running an equivalent SCALE/TRITON model on the same 2D slice and quantifying k\_eff bias as a function of energy-group structure.
\item How does the reactivity evolution change under a 3D axial model with realistic axial power profile and graphite growth (anisotropic swelling) over 2+ years
\end{enumerate}
\paragraph{Variational Monte Carlo calculations of A\textless{}=4 nuclei with an artificial neural-network correlato \textnormal{\small(OSTI \texttt{1864334})}}
\begin{enumerate}[leftmargin=1.6em,itemsep=0.2em,topsep=0.2em,label=Q\arabic*.]
\item Does adding a differentiable backflow transformation (learned per-pair coordinate shifts) to the symmetric ansatz close most of the 6 MeV Minnesota-central 4He gap, and how does the closure scale with ansatz width
\item For A=4 under Minnesota, how does the stochastic-reconfiguration optimiser compare to Adam in (a) final energy and (b) sample efficiency at matched wall time on a single A100
\item Can a permutation-equivariant transformer correlator generalise across A=2..4 in a single training run (transfer-learning the pair correlator with additional three-body heads)
\item What is the minimum network capacity (parameters, depth, width) at which the deuteron accuracy saturates at the exact Minnesota bound, and does that minimum scale super-linearly with A
\item When the Minnesota potential is replaced by a spin-orbit-augmented version (Minnesota + LS), does the NN correlator recover the spin-orbit-induced tensor correlations, or must they be supplied via an operator ansatz
\end{enumerate}
\subsection*{Particle Physics}
\addcontentsline{toc}{subsection}{Particle Physics}
\paragraph{Spin-Dependent Scattering of Sub-GeV Dark Matter \textnormal{\small(OSTI \texttt{3014512})}}
\begin{enumerate}[leftmargin=1.6em,itemsep=0.2em,topsep=0.2em,label=Q\arabic*.]
\item How does the SD reach for Al2O3 change if realistic experimental resolution (sigma\_E \textasciitilde{} 0.5 meV TES) and backgrounds are folded into the 1 meV threshold curves
\item Run the DarkELF computation on a wider target panel (sapphire, GaN, diamond, ZnO, Si3N4) at 1 meV threshold and rank candidates by SD reach per kg$\cdot$yr for each operator.
\item Produce the missing Figs. 9-10 by scanning m\_med over [0.1, 10]*q\_0 per m\_chi and constructing the envelope minimum-sigma curve; quantify how much reach improves vs fixed m\_med=3q\_0.
\item For the A' model, fold in the realistic SN+SIDM constraint band (from the paper's Figs. 1-4) and identify the residual DM-mass window where Al2O3 at 1 meV can make a world-leading statement.
\item Compare the DarkELF single-phonon vs multi-phonon+impulse transition by checking the crossover energy and cross-section continuity at E\_th \textasciitilde{} phonon bandwidth for Al2O3.
\end{enumerate}
\subsection*{Power/Electronics}
\addcontentsline{toc}{subsection}{Power/Electronics}
\paragraph{Common Mode Voltage Reduction of Single-Phase Quasi-Z-Source Inverter-Based Photovoltaic System \textnormal{\small(OSTI \texttt{1606674})}}
\begin{enumerate}[leftmargin=1.6em,itemsep=0.2em,topsep=0.2em,label=Q\arabic*.]
\item How does the 50\% CMV reduction degrade as the CM parasitic capacitance C\_g varies over 0.5-10 nF (representative of different PV-array sizes and mounting conditions)
\item Does the ZS1-\textgreater{}ZS2 modulation preserve its CMV advantage under dead-time insertion of 0.5-2 us, or does dead-time reintroduce intermediate CM levels
\item Can a third modulation state (ST-\textgreater{}ZS3) further reduce CMV peak below V\_PN/4 without penalising DC-link utilisation, and what is the achievable theoretical floor
\item How do the CMV and leakage-current reductions translate to measured conducted EMI spectra (150 kHz-30 MHz) on a prototype, and does a smaller EMI filter become viable
\item When the grid voltage is distorted (\textgreater{}5\% THD, notched by nonlinear loads), does the proposed scheme maintain the clean two-level CMV, or does the control interaction reintroduce transients
\end{enumerate}
\paragraph{Learning Sequential Distribution System Restoration via Graph-Reinforcement Learning \textnormal{\small(OSTI \texttt{1868518})}}
\begin{enumerate}[leftmargin=1.6em,itemsep=0.2em,topsep=0.2em,label=Q\arabic*.]
\item Run the trained GCN-DQN on a full IEEE 8500-node case (e.g. from OpenDSS test feeders) and measure zero-shot transfer vs fine-tuned performance --- does the GCN's topology-aware embedding genuinely generalize
\item Compare GCN-DQN against modern graph-transformer policies (e.g. Graphormer, GPS) with identical training budgets to test whether the paper's GCN advantage holds under stronger architectural baselines.
\item Formulate the restoration problem as MILP in pandapower+pyomo and quantify the optimality gap of GCN-DQN --- how close does RL get to CPLEX-optimal on networks where MILP is tractable
\item Add adversarial fault scenarios (targeted attacks rather than random 10-30\% line outages) and measure GCN-DQN robustness vs the greedy heuristic baseline.
\item Study reward-shaping alternatives: compare the paper's weighted load/voltage/thermal reward against a strict load-restoration-only reward and a constraint-violation-barrier reward --- which gives the most reliable policy in realistic fault scenarios
\end{enumerate}
\subsection*{Quantum/AMO}
\addcontentsline{toc}{subsection}{Quantum/AMO}
\paragraph{Exactly-Solved Model of Light-Scattering Errors in Quantum Simulations with Metastable Trapped- \textnormal{\small(OSTI \texttt{2441075})}}
\begin{enumerate}[leftmargin=1.6em,itemsep=0.2em,topsep=0.2em,label=Q\arabic*.]
\item Implement the full 2D Penning-trap equilibrium-crystal mode solver and recompute J\_ij --- do the Fig 2 fidelity crossovers match the paper's reported N\_threshold at B=0.9T and B=4.5T within 10\%
\item Extend the exact solution from $\pi$ to $\sigma$$\pm$ polarization where |1 is no longer dark --- derive new analytic expressions and quantify the fidelity penalty of losing the dark-state protection.
\item Test the exact solution against a full Lindblad simulation at N=20 (near the boundary of tractable numerics) --- does the product-form correlation factorization still agree to machine precision, or do many-body corrections emerge
\item For $^3$Ba (different branching ratios, \textasciitilde{}70\% leakage), do the novel L and B terms still dominate the error budget, and does postselection recover the same 5.5\% residual-error advantage
\item Combine the exact scattering model with realistic single-qubit gate errors (miscalibration, Stark shifts) to produce a full end-to-end error budget for a 50-qubit GHZ preparation circuit --- at what N does gate error overtake scattering as the dominant error source
\end{enumerate}
\paragraph{FDTD: Solving 1+1D Delay PDE in Parallel \textnormal{\small(OSTI \texttt{1475143})}}
\begin{enumerate}[leftmargin=1.6em,itemsep=0.2em,topsep=0.2em,label=Q\arabic*.]
\item Port the solver to CUDA and measure sustained throughput for (N, N) up to the Fig 8 regime --- at what grid does the GPU cross the memory wall, and does out-of-core NVMe backing restore feasibility
\item Implement the exact scattering-theory reference (Ref [8]) and quantify FDTD g error as a function of $\Delta$ --- does the empirical convergence rate match the paper's theoretical O($\Delta$$^2$)
\item Extend the PDE to N atoms at staggered positions --- does g(0) show monotonic deepening of antibunching, and is there a critical atom spacing where superradiance flips to subradiance
\item How does g(0) depend on ka beyond $\pi$/2 up to the Fig 8 regime (ka $\approx$ 10$\pi$--20$\pi$) Produce a phase diagram of g(0) vs (ka, $\omega$/) showing antibunching/bunching/chaotic zones.
\item Replace the perfect mirror with a finite-reflectivity boundary r\textless{}1 --- at what r does the non-Markovian oscillation structure collapse to Markovian exponential decay, and does an effective-Markovian approximation work beyond that threshold
\end{enumerate}
\paragraph{Simple Coplanar Waveguide Resonator Mask Targeting Multiplexed Superconducting-Qubit Readout \textnormal{\small(OSTI \texttt{1983793})}}
\begin{enumerate}[leftmargin=1.6em,itemsep=0.2em,topsep=0.2em,label=Q\arabic*.]
\item Run the same participation analysis in Palace (AWS open-source FEM) on the published GDS mask --- do the absolute p\_MS and p\_MA values drop toward the paper's $\times$10$^3$ scale when sub-nm interface meshing is used
\item Sweep trench depth 0--300 nm in 10 nm steps --- does p\_SA decrease monotonically, and where is the diminishing-returns knee for TLS loss reduction
\item How sensitive is p\_MS/p\_SA $\approx$ 2 to center-conductor-to-gap aspect ratio s/g Sweep s $\in$ \{3,6,9,12\} m holding Z=50 $\Omega$ via g adjustment; report the ratio's $\pm$3$\sigma$ band across realistic fab tolerances.
\item Model kinetic inductance of a 100 nm Nb film with $\lambda$\_L = 80 nm and re-compute Z and f --- does the 3.4\% gap close, and do all 8 frequencies shift consistently
\item Couple the analytical Q\_c(\_c) curve to the participation-ratio loss model --- can a single scalar 'design efficiency' metric recommend an optimal coupler length that jointly minimizes TLS loss and maximizes multiplexed readout fidelity
\end{enumerate}

\subsection{Wave 3 (2026-05-05): Biology and BV-BRC Batch}
Wave~3 added 10 biology/bioinformatics replications in a single day. Key
advances:

\begin{itemize}
\item \textbf{BV-BRC API as primary data source.} Five papers used the BV-BRC
(Bacterial and Viral Bioinformatics Resource Center) API for genome retrieval,
AMR detection, virulence gene profiling, MLST typing, and pan-genome
analysis. The API proved highly effective: all 5 papers achieved $\geq$67\%
claim coverage with 0\% contradiction rate.
\item \textbf{Sequence-level bioinformatics.} Three papers (Jiang 2017 ARG
dissemination, Price 2018 RB-TnSeq, Vassallo 2022 phage defence) involved
direct sequence analysis (BLASTP, EMBOSS needle, bowtie2). These achieved
the highest precision in the biology cohort: Jiang 2017 matched 56/56 ARG
protein identities with mean deviation of 0.3\%.
\item \textbf{Largest-scale single-paper replication.} Price et al.\ 2018
covered 32 organisms, 150,851 protein-coding genes, and 4,870 experiments---
the most complex biology replication in the project and one of the
largest-scale replications overall.
\item \textbf{Spot-check limitations.} Sato et al.\ 2021 (centenarian bile
acids) scored lowest (4/6) due to read subsampling constraints (3--7\% of
total depth). This underscores that shallow metagenomics spot-checks have
limited discriminative power for rare-organism enrichment claims.
\item \textbf{Tool substitution documentation.} BV-BRC PLFam for Roary,
abricate for CARD/RGI, FastANI for pyani---all produce consistent
qualitative conclusions despite quantitative differences in absolute counts.
These substitutions are documented and defensible.
\end{itemize}

\noindent\textbf{Domain coverage.} The 10 biology papers span 7 sub-domains:
microbial AMR epidemiology (Jiang, Zhang), functional genomics (Price),
phage defence (Vassallo), microbiome ecology (Li), metagenomics (Sato),
taxonomic genomics (Fluit), veterinary genomics (Sivakumar), metabolic
annotation (Shrestha), and bacterial pan-genomics (Thakur). This is the
first systematic biology batch and establishes the BV-BRC API as a viable
replication platform.

\noindent\textbf{Wave~3 summary statistics (10 papers):}
mean coverage = 7.6/10, mean agreement = 8.1/10.
7/10 scored REPLICATED, 3/10 PARTIAL or SPOT-CHECK.
0/10 contradicted.

\subsection{Wave 4 (2026-05-07): PDE and PINN Batch}
Wave~4 added 5 PDE/PINN replications in a single day, probing the
reproducibility of machine-learning approaches to partial differential
equations. Key findings:

\begin{itemize}
\item \textbf{Code availability as the decisive factor.} The two papers with
  public GitHub repositories (Kochkov~2021 jax-cfd, Grossmann~2023
  FEM-vs-PINNs) both achieved REPLICATED status (C=8, A=9 each). The three
  papers without public code (Eivazi~2022, Kopani\'cakov\'a~2023,
  Zhang~2019) all achieved PARTIAL, with absolute errors 10--100$\times$
  higher than claimed.
\item \textbf{PINN reproducibility gap.} The three PARTIAL papers share a
  common failure pattern: qualitative methodology reproduces correctly, but
  quantitative precision does not. This is attributable to: (i)~missing
  custom L-BFGS optimizers, (ii)~undocumented hyperparameters, and
  (iii)~unavailable reference datasets. See \S\ref{sec:pinn-repro} for
  detailed analysis.
\item \textbf{Data-blocker resolution.} Wave~4 demonstrated two creative
  data-blocker workarounds: self-generating ground truth via analytical
  solutions (3D Poisson) and split-step Fourier methods (1D Schr\"odinger),
  upgrading Grossmann~2023 from PARTIAL to REPLICATED.
\item \textbf{Author model verification.} For Kochkov~2021, we independently
  evaluated the authors' pre-trained model outputs (released on GCS),
  providing a gold-standard comparison separate from our own training.
\end{itemize}

\noindent\textbf{Wave~4 summary statistics (5 papers):}
mean coverage = 7.2/10, mean agreement = 7.0/10.
2/5 scored REPLICATED, 3/5 PARTIAL.
0/5 contradicted.

